

\chapter*{ABSTRACT}
\phantomsection
\addcontentsline{toc}{chapter}{ABSTRACT}


\bigskip
Finding a person's location in remote mountain regions can be difficult when satellite navigation is unavailable or unreliable. In deep valleys, surrounding mountains can block \gls{gnss} signals, while mobile network coverage is often limited. This creates a need for an alternative method that can estimate location without depending on satellite positioning or internet connectivity. This project presents an offline visual geolocation system for the Khumbu region of Nepal that estimates the location of a photograph using the surrounding mountain skyline. A \gls{dem} is used to generate a database of reference skylines from regularly spaced viewpoints across the study area. When a user captures a photograph, the skyline is extracted by separating the sky from the terrain using a lightweight semantic segmentation model based on the U-Net architecture with a MobileNetV3 encoder. Because labelled Himalayan images are scarce, the segmentation model is trained on GeoPose3K's predefined training split (3,000 real images) combined with 300 synthetically generated images created from the terrain model under different illumination and sky appearance conditions. The extracted skyline is then compared with the reference database using a two-stage hierarchical matching process. A fast coarse search identifies the most likely candidates, followed by a finer alignment stage that improves robustness against small viewpoint and camera differences. The proposed method was evaluated using 300 synthetic test images covering different locations and environmental conditions within the study area. Under grid-aligned synthetic evaluation with optional height-filtered candidate pruning, where query viewpoints coincide with reference grid locations and can be filtered using approximate height information, the system correctly localized 94.67\% of the test images within 1 km of the true location and 84.00\% within 100 m. These results suggest that mountain skylines provide reliable geographic information for visual localization in the Khumbu region. Although the current evaluation is based on synthetic imagery, the framework establishes a practical foundation for lightweight offline mountain navigation and can be extended through future testing with real-world photographs.

\vspace{1em}
\noindent\textit{\textbf{Keywords}}--- visual geolocation, mountain skyline, \acrshort{dem}, semantic segmentation, synthetic data, Himalayan navigation, offline localization
